\section{Experimental Assessment of the Built Cooker}
\label{sec:experimental_assessment}

After construction, the simulated optical performance should be checked with a small set of
repeatable outdoor tests.  The measurements should follow the spirit of ASAE S580.1 solar
cooker testing, as used in Solar Cookers International's Performance Evaluation Process,
which reports water-heating power from measured temperature rise, solar irradiance, ambient
temperature, and wind speed.\footnote{\url{https://www.solarcookers.org/files/7814/2422/1265/ASAE-standard-s-5801-aug13.pdf};
\url{https://www.solarcookers.org/resources/results}.}  Each test should record the date,
location, ambient temperature, direct solar irradiance, wind speed, tracking interval, lid
configuration, and initial load temperature.

\subsection{Stagnation Temperature}

The first test is a no-load stagnation run.  The cooker is aligned toward the sun with the
chosen lid installed and no water or cooking load inside.  Thermocouples are placed on the
matte bottom plate, on the thermal battery top, and on the inner air space.  The maximum
steady temperature indicates the upper thermal potential of the design and can be reported
with the box-cooker first figure of merit,
\begin{equation}
  F_1 = \frac{T_\mathrm{stag} - T_a}{G},
\end{equation}
where $T_\mathrm{stag}$ is the stagnation temperature, $T_a$ is ambient temperature, and
$G$ is the measured solar irradiance.

\subsection{Water Heating and Boiling}

The main performance test uses a blackened vessel containing a known mass of water.  A
standard practical load for this prototype is $1.0$ L of water, although an ASAE-style load
scaled by cooker intercept area may also be reported.  Water temperature is logged every
5--10 min until boiling or until the temperature plateaus.  The useful heating power over each
time interval is
\begin{equation}
  P_i = \frac{m_w c_p (T_{w,2}-T_{w,1})}{\Delta t},
\end{equation}
with $c_p = 4186$ J kg$^{-1}$ K$^{-1}$.  To compare tests taken under different sunlight,
the interval power is normalized to 700 W m$^{-2}$,
\begin{equation}
  P_s = P_i \frac{700}{G_i}.
\end{equation}
The report should include the time to reach 60, 80, 95, and 100 $^\circ$C, the peak water
temperature, the average standardized cooking power, and whether a rolling boil was reached.

\subsection{Design Comparison Tests}

The same water-heating test should be repeated for the main build variants: single-pane PMMA,
double-pane PMMA, and, if available, equivalent glass lids.  This directly measures the
trade-off between optical transmission losses and reduced thermal losses through the opening.
The thermal battery should also be instrumented during these runs, with sensors on the top and
four side walls, so the measured heating pattern can be compared against the ray-tracing flux
maps.

\subsection{Retention and Usability}

After a heating run, the cooker should be shaded or closed and allowed to cool while water,
battery, bottom-plate, and air temperatures are logged.  This cooldown curve estimates heat
retention and post-sun cooking usefulness.  A final practical check should cook a simple
repeatable food load, such as rice or potatoes, using the same tracking schedule expected in
normal use.  Results should be summarized in a compact table containing the lid type,
irradiance range, wind speed, peak temperatures, time to boil 1 L, standardized cooking power,
and cooldown time from 90 to 60 $^\circ$C.
